\section{Testen des Algorithmus}\label{sec:testen_des_algorithmus}
In diesem Unterkapitel werden die Test-Szenarien auf die der Tester getestet wird beschreiben. Außerdem wird die Beschaffung der Testdaten erläutert und die Testergebnisse bewertet. 

\subsection{Allgemeine Bewertung}\label{subsec:allgemeine_bewertung_testen}
Der Algorithmus ist bereits in der Lage Gipfel zu erkennen. Er hat jedoch Probleme mit den Randwerten und der Dominanz sowie Prominenzberechnung. Daher müssen diese Programmbausteine untersucht und verbessert werden. Besonders auffällig sind auch die Hohen Höhenabweichungen bei der Wetterstein und der Südschwarzwaldkarte. Um diese bei künftigen Tests als Fehler auszuschliesen sollten die Geodaten des Copernicus-Programm nicht mehr genutzt werden.
